\documentclass[conference]{IEEEtran}

\usepackage[T1]{fontenc}
\usepackage[utf8]{inputenc}
\usepackage[spanish,es-tabla]{babel}
\usepackage{lmodern}
\usepackage{microtype}
\usepackage{csquotes}
\usepackage{amsmath,amssymb}
\usepackage{booktabs}
\usepackage{tabularx}
\usepackage{array}
\usepackage{url}
\usepackage[hidelinks]{hyperref}
\usepackage[backend=biber,style=ieee,sorting=none]{biblatex}
\addbibresource{referencias.bib}

\usepackage{enumitem}
\setlist[enumerate]{leftmargin=*,itemsep=1pt,topsep=2pt}
\setlist[itemize]{leftmargin=*,itemsep=1pt,topsep=2pt}
\setlength{\parindent}{0.65em}
\setlength{\parskip}{1pt}
\setlength{\columnsep}{0.24in}
\renewcommand{\arraystretch}{1.16}
\newcolumntype{Y}{>{\raggedright\arraybackslash}X}
\newcommand{\req}[1]{\noindent\textbf{#1:}\enspace}
\DefineBibliographyStrings{spanish}{bibliography={Referencias}}
\urlstyle{same}

\title{Comparativa de Modelos de Proceso de Requisitos en Escenarios Reales}

\author{
	\IEEEauthorblockN{José Ignacio Gilces Carranza}
	\IEEEauthorblockA{Universidad Técnica Estatal de Quevedo (UTEQ) --- Facultad de Ciencias de la Computación\\
		Ingeniería de Software $\cdot$ Ingeniería de Requerimientos (ISR-401) $\cdot$ Práctica en clase FORM-PC-01}
}

\begin{document}
	\onecolumn
	\maketitle
	
	\begin{abstract}
		En este informe se formaliza el desarrollo preliminar de la actividad \texttt{FORM-PC-01}, cuyo objetivo es comparar los modelos de proceso de ingeniería de requisitos (cascada/secuencial, iterativo/incremental y ágil/Scrum) en cuatro escenarios representativos del entorno profesional. Se definen cinco criterios cuantitativos y cualitativos de selección y se formula la matriz comparativa correspondiente.
	\end{abstract}
	
	\section{Criterios de Comparación de Modelos de Proceso}
	La selección de un modelo de proceso de requisitos no responde a decisiones arbitrarias del equipo de desarrollo, sino a un análisis riguroso del dominio del problema y del entorno operativo \cite{pohl2025, swebok2024}. Se establecen cinco criterios de evaluación:
	
	\begin{enumerate}
		\item \textbf{C1. Estabilidad y volatilidad de los requisitos:} Mide la frecuencia y magnitud de las modificaciones en las necesidades del sistema durante el ciclo de vida. Una alta volatilidad favorece enfoques adaptativos; una alta estabilidad permite planificaciones secuenciales fijas.
		\item \textbf{C2. Criticidad del sistema y nivel de riesgo:} Evalúa el impacto potencial de fallos en términos de seguridad humana (\emph{safety}), seguridad de la información (\emph{security}) o pérdidas financieras. A mayor criticidad, mayor rigidez en la verificación temprana.
		\item \textbf{C3. Disponibilidad e involucramiento del cliente/stakeholder:} Cuantifica la capacidad y disposición de las partes interesadas para participar activamente en sesiones continuas de elicitación, retroalimentación y validación de incrementos.
		\item \textbf{C4. Rigidez contractual y regulatoria:} Determina la presencia de marcos legales, leyes nacionales (como la LOPDP \cite{lopdp2021}), plazos fijos de contratación pública o normativas institucionales que imponen un alcance cerrado antes del inicio del desarrollo.
		\item \textbf{C5. Tiempo para el primer incremento funcional (\emph{Time-to-Market}):} Mide la urgencia del negocio por desplegar una versión operativa en producción para capturar usuarios, validar modelo comercial o mitigar riesgos de mercado.
	\end{enumerate}
	
	\section{Matriz Comparativa de Escenarios y Modelos}
	Se evalúan los cuatro escenarios del caso frente a los tres modelos principales de proceso de requisitos.
	
	\begin{center}
		\centering
		\scriptsize
		\setlength{\tabcolsep}{3pt}
		\begin{tabularx}{\textwidth}{@{}p{2.6cm} Y p{1.8cm} p{1.8cm} p{1.8cm} p{2.2cm}@{}}
			\toprule
			\textbf{Escenario Evaluado} & \textbf{Características del Entorno} & \textbf{Modelo Cascada (Secuencial)} & \textbf{Modelo Iterativo / Incremental} & \textbf{Modelo Ágil (Scrum)} & \textbf{Modelo Seleccionado}\\
			\midrule
			\textbf{Escenario 1:} Sistema embebido de planta industrial. & Alta criticidad de hardware, tolerancia cero a fallos, requisitos de control físico estables. & \textbf{Alta} (Fase de requisitos cerrada antes de fabricar). & Media (Limitada por despliegue físico). & Baja (Dificultad de refactorizar firmware físico). & \textbf{Modelo Cascada (Secuencial)} \\
			\addlinespace[3pt]
			\textbf{Escenario 2:} Aplicación móvil de emprendimiento. & Requisitos altamente volátiles, necesidad de probar mercado móvil e iterar rápido. & Baja (Riesgo de obsolescencia antes de lanzar). & Media (Ciclos largos para startups). & \textbf{Alta} (Sprints rápidos con retroalimentación). & \textbf{Modelo Ágil (Scrum)} \\
			\addlinespace[3pt]
			\textbf{Escenario 3:} Sistema de un GAD con contrato público. & Alcance legalmente cerrado por plazos, pero alta necesidad de validar módulos operativos. & Media (Riesgo de rechazo en entrega final). & \textbf{Alta} (Control por hitos y entregables). & Baja (Dificultad de contratos fijos por Sprints). & \textbf{Modelo Iterativo / Incremental} \\
			\addlinespace[3pt]
			\textbf{Escenario 4:} Plataforma de cooperativa de ahorro. & Regulación financiera estricta, alta seguridad, alta concurrencia y despliegue por módulos. & Media (Retrasa validación de transacciones). & \textbf{Alta} (Incrementos modulares y seguros). & Media (Requiere gobierno estricto de cambios). & \textbf{Modelo Iterativo / Incremental} \\
			\bottomrule
		\end{tabularx}
	\end{center}
	
	\section{Referencias}
	\printbibliography[heading=none]
	
\end{document}